\section{Riesgo de complejidad: me preocupa que Manus se vuelva complejo}

Las secciones anteriores trataron la entrega al estilo de la manufactura y la infraestructura; esta sección vuelve a la filosofía de producto. Al final de la entrevista, Peak dijo que lo que más teme es que Manus pierda su carácter distintivo, y que lo que más se teme dentro del equipo es que Manus se vuelva complejo. Este juicio es muy importante. Muchos productos, al crecer, agregan funciones de manera natural para atender a más usuarios, más escenarios y más clientes; pero cada cosa que se agrega diluye a las demás.

\begin{figure}[H]
\centering
\includegraphics[width=0.96\textwidth]{complexity-risk.png}
\caption{Riesgo de complejidad: crecimiento, funciones, dilución, contención y volante de inercia.}
\end{figure}

\begin{knowledgebox}{Lectura de la figura: la complejidad nace del atractivo del crecimiento}
La figura avanza del crecimiento a las funciones y, después, a la dilución y la contención. El crecimiento del producto lleva al equipo a añadir más capacidades, pero cuantas más funciones hay, más le cuesta al usuario entender el valor central. El reto de Manus es mantener la contención sin sacrificar el crecimiento sostenido.
\end{knowledgebox}

\subsection{No se teme a la competencia, sino a perder el valor distintivo}

Peak considera que competir con las grandes empresas no es en sí lo más temible; cuando te preocupa que te corten el suministro o que te copien, el producto ya llegó a un muy buen estado. El riesgo mayor es perder el valor exclusivo, que los usuarios ya no sepan por qué elegirte. Esto también es una forma de competencia, pero no es la competencia externa la que te elimina primero, sino la complejidad interna la que primero te diluye.

\begin{importantbox}{El límite mínimo de la singularidad de un producto}
Un producto puede ampliarse, pero no puede hacer que el usuario olvide cuál es el problema central que resuelve. El crecimiento no puede tener como costo perder el lugar que el producto ocupa en la mente del usuario.
\end{importantbox}

\subsection{Resumen de la sección}

La ansiedad de Manus ante la complejidad es un problema que enfrentarán todas las empresas de Agent: cuanto más éxito tienen, más fácil es agregar funciones, y cuantas más funciones agregan, más probable es que pierdan un lugar claro en la mente del usuario.
